Read in Gauge Theory Gravity: the metric $g$ below is the derived metric of the gravitational field $(\hb, \Omega)$ on flat spacetime, the manifold is the invertible region, and the tractor connection is the flat-space $\Cl(2,4)$ gauge field of Theorem \claimref{synthesis.gtg-tractor-gauge-field-reproduces-tractor-connection}; the index computation below is its reading in the splitting of $g$ (Fact \claimref{synthesis.imported-gtg-dictionary}).

(1) Boost part: $\tau\sigma_i = p_i + O(\tau^{2/3})$ by Lemma \claimref{interior.exact-parallelism-convergence-rate}(2), continuous up to $\tau = 0$ with boundary value $p_i$; in the admissible class continuity on unit directions is exactly the hypothesis $E_\tau \to 0$ of Definition \claimref{obj.interior.weighted-blow-up-hubble-normalisation-admissible}. Rotation part: by Lemma \claimref{interior.interior-connection-exactly} the rotation coefficients on unit directions are $O(1/r)$, so Hubble-normalised they are $O(\tau/r) = O(\tau^{1/3}) \to 0$. Tractor connection: the Einstein hypothesis of Definition \claimref{obj.interior.weighted-blow-up-hubble-normalisation-admissible} gives $\Sch_{ab} = \tfrac{\Lambda}{6}\,g_{ab}$ in the scale $g$ (Lemma \claimref{interior.metric-parallelism-trace-schouten-condition}(ii), Theorem \claimref{interior.scale-tractor-leaves-cone}), so by \eqref{eq:tractorconn} the tractor connection is the Levi-Civita action on the $\mu$-slot plus coupling maps ($\mu \mapsto -\mu_a$ into the $\sigma$-slot, $\rho \mapsto g_{ab}\rho$ and $\sigma \mapsto -\tfrac{\Lambda}{6}g_{ab}\sigma$ into the $\mu$-slot, $\mu \mapsto \tfrac{\Lambda}{6}\mu_a$ into the $\rho$-slot) whose coefficients on unit directions are $\tau$-independent, so their Hubble normalisation is $O(\tau) \to 0$; without a field-equation hypothesis the Schouten term is not controlled relative to $g$ and no such statement holds, which is why the hypothesis is part of admissibility. (2) $I_g$ is parallel exactly, at every $\tau > 0$ (Theorem \claimref{interior.covariant-form-vacuum-equations} for $\Lambda = 0$, Theorem \claimref{interior.scale-tractor-leaves-cone} in general), and in the parallel tractor frame it is a constant vector, so it extends over the boundary as the constant section of the boundary bundle. Parallelism at the boundary is pointwise algebra: the limiting connection values are the weighted boosts $-p_i B_i$ by (3), each $B_i = e_0 \wedge e_i$ has both factors in the $\mu$-slot subspace $W$, and the tractor metric pairs $W$ to zero against $I_g$, whose $\mu$-component vanishes; so by \eqref{eq:bivectoraction}, $B_i(I_g) = \ip{e_i}{I_g}\,e_0 - \ip{e_0}{I_g}\,e_i = 0$, and a constant section annihilated by every value of the limiting connection is parallel for it. (3) The surviving matrix is read off (1): boost coefficients $-\tau\sigma_i$ converge to $-p_i$, rotations and couplings to zero. (4) Containment: the limiting values lie in $\operatorname{span}(B_1, B_2, B_3) \subset \Lambda^2 W$. Fullness: every $p_i$ is nonzero, so each of the three boosts is itself a limiting value, and the Lie algebra the three generate is exactly $\Lambda^2 W \cong \mathfrak{so}(1,3)$ by Lemma \claimref{interior.boosts-generate-exactly-levi-factor}; the boundary holonomy algebra of Definition \claimref{obj.interior.weighted-blow-up-hubble-normalisation-admissible}, the algebra generated by the values, is therefore exactly $\Lambda^2 W$. In the decomposition of Proposition \claimref{interior.stabilisers-flat-model}, the two nilpotent summands $\nInf\wedge W$, $\nO\wedge W$ are the couplings that scale to zero, and the grading element $\nInf \wedge \nO$ receives no connection because the $\sigma$- and $\rho$-slots carry no diagonal term in \eqref{eq:tractorconn}, both being fed from the $\mu$-slot alone whatever the Schouten tensor. The reading of $\Lambda^2 W$ inside the stabiliser depends on $\Lambda$. For $\Lambda = 0$ the frozen point is the null $I_g = (1,0,0)$ and $\Lambda^2 W$ is precisely the Levi factor of its parabolic stabiliser. For $\Lambda \neq 0$ the frozen point $(1, 0, \Lambda/6)$ has norm $\Lambda/3 \neq 0$, so by Proposition \claimref{interior.non-null-stabilisers-section-they-act}(1) its stabiliser is the reductive $\Lambda^2(I_g^\perp)$, and since $W \perp I_g$ the six-dimensional $\Lambda^2 W$ lies inside that ten-dimensional algebra as a proper subalgebra. The computed algebra is the same for every $\Lambda$; only its identification as a Levi factor is a $\Lambda = 0$ statement. Final clause: by Lemma \claimref{interior.interior-connection-exactly} the connection along the $t$-lines is $-\sigma_1 B_1\,\theta^1$ and along the $\theta$-lines is $-\sigma_2 B_2\,\theta^2$, in each case a fixed bivector times a scalar form, so all transports in one such direction commute. Hubble-normalised, the interior generator on the unit direction is $-(\tau\sigma_i)B_i$ against the model's $-p_iB_i$, and $\tau\sigma_i - p_i = O(\tau^{2/3})$ by Lemma \claimref{interior.exact-parallelism-convergence-rate}(2): over segments of fixed Hubble-normalised length the transports differ by a boost of angle $O(\tau^{2/3}) \to 0$. Over a fixed \emph{coordinate} segment the generator is $-(\partial_\tau a_i)B_i$, of size $\tau^{p_i - 1}$, the transports individually diverge, and no claim is made about their comparison. The frame comparison along the $\tau$-geodesics is a separate scalar computation: $\int_0^\tau(\sigma_i - p_i/s)\,ds = \log\bigl(a_i(\tau)\,\tau^{-p_i}/c_i\bigr)$ with $c_i := \lim_{s \to 0} a_i(s)\,s^{-p_i}$, finite and nonzero by Lemma \claimref{interior.kasner-form-interior}, so the coordinate frame's evolution converges to the model's at rate $O(\tau^{2/3})$, direction by direction, with no noncommutativity anywhere in the argument. Along the $\varphi$-lines the connection on the coordinate vector $\partial_\varphi$ is $(\partial_\tau a_3)\,B_3 + \cos\theta\,R_{23}$ (Lemma \claimref{interior.interior-connection-exactly}), a divergent boost beside a $\tau$-independent rotation, and the two do not commute: a Duhamel comparison would conjugate the $O(\tau^{1/3})$ discrepancy by an unbounded model transport, which the integrability hypothesis alone does not control, so no transport-level claim is made in that direction here; Theorem \claimref{interior.transport-level-convergence-along-tangential-curves} and Corollary \claimref{interior.schwarzschild-lines-exactly} below make it, for segments of fixed Hubble-normalised length, where the conjugation is harmless. Clause (4) asserts the values-level algebra only, and that is what every later use in the series consumes.
